% TODO proof-reading
\subsubsection{可変引数関数}

可変引数関数は実際には配列を使用します：


\begin{lstlisting}[style=customjava]
	public static void f(int... values)
	{
		for (int i=0; i<values.length; i++)
			System.out.println(values[i]);
	}

	public static void main(String[] args) 
	{
		f (1,2,3,4,5);
	}
\end{lstlisting}

\begin{lstlisting}
  public static void f(int...);
    flags: ACC_PUBLIC, ACC_STATIC, ACC_VARARGS
    Code:
      stack=3, locals=2, args_size=1
         0: iconst_0      
         1: istore_1      
         2: iload_1       
         3: aload_0       
         4: arraylength   
         5: if_icmpge     23
         8: getstatic     #2       // Field java/lang/System.out:Ljava/io/PrintStream;
        11: aload_0       
        12: iload_1       
        13: iaload        
        14: invokevirtual #3       // Method java/io/PrintStream.println:(I)V
        17: iinc          1, 1
        20: goto          2
        23: return        
\end{lstlisting}

\ttfはオフセット3の\TT{aload\_0}を使用して整数の配列を受け取ります。

そして配列のサイズを取得します、等々。


\begin{lstlisting}
  public static void main(java.lang.String[]);
    flags: ACC_PUBLIC, ACC_STATIC
    Code:
      stack=4, locals=1, args_size=1
         0: iconst_5      
         1: newarray       int
         3: dup           
         4: iconst_0      
         5: iconst_1      
         6: iastore       
         7: dup           
         8: iconst_1      
         9: iconst_2      
        10: iastore       
        11: dup           
        12: iconst_2      
        13: iconst_3      
        14: iastore       
        15: dup           
        16: iconst_3      
        17: iconst_4      
        18: iastore       
        19: dup           
        20: iconst_4      
        21: iconst_5      
        22: iastore       
        23: invokestatic  #4       // Method f:([I)V
        26: return        
\end{lstlisting}

配列は\main内で\TT{newarray}命令を使用して構築され、それから埋められ、\ttfが呼び出されます。


ところで、配列オブジェクトは\mainの終わりで破棄されません。

\myindex{Garbage collector}
Javaにはデストラクタが全くありません。なぜならJVMにはガベージコレクタがあり、必要と感じたときに自動的にこれを行うからです。


\TT{format()}メソッドはどうでしょうか？

入力として2つの引数を取ります：文字列とオブジェクトの配列：


\begin{lstlisting}
	public PrintStream format(String format, Object... args)
\end{lstlisting}
( \url{http://docs.oracle.com/javase/tutorial/java/data/numberformat.html} )

見てみましょう：

\begin{lstlisting}[style=customjava]
	public static void main(String[] args)
	{
		int i=123;
		double d=123.456;
		System.out.format("int: %d double: %f.%n", i, d);
	}
\end{lstlisting}

\begin{lstlisting}
  public static void main(java.lang.String[]);
    flags: ACC_PUBLIC, ACC_STATIC
    Code:
      stack=7, locals=4, args_size=1
         0: bipush        123
         2: istore_1      
         3: ldc2_w        #2         // double 123.456d
         6: dstore_2      
         7: getstatic     #4         // Field java/lang/System.out:Ljava/io/PrintStream;
        10: ldc           #5         // String int: %d double: %f.%n
        12: iconst_2      
        13: anewarray     #6         // class java/lang/Object
        16: dup           
        17: iconst_0      
        18: iload_1       
        19: invokestatic  #7         // Method java/lang/Integer.valueOf:(I)Ljava/lang/Integer;
        22: aastore       
        23: dup           
        24: iconst_1      
        25: dload_2       
        26: invokestatic  #8         // Method java/lang/Double.valueOf:(D)Ljava/lang/Double;
        29: aastore       
        30: invokevirtual #9         // Method java/io/PrintStream.format:(Ljava/lang/String;[Ljava/lang/Object;)Ljava/io/PrintStream;
        33: pop           
        34: return        
\end{lstlisting}

従って\IT{int}と\IT{double}型の値は、まず\TT{valueOf}メソッドを使用して\TT{Integer}と\TT{Double}オブジェクトに昇格されます。

\TT{format()}メソッドは入力として\TT{Object}型のオブジェクトを必要とし、\TT{Integer}と\TT{Double}クラスはルート\TT{Object}クラスから派生しているため、入力配列の要素として適しています。

一方、配列は常に同質的、つまり異なる型の要素を保持できないため、\IT{int}と\IT{double}値をそれに押し込むことは不可能です。


\TT{Object}オブジェクトの配列がオフセット13で作成され、\TT{Integer}オブジェクトがオフセット22で配列に追加され、\TT{Double}オブジェクトがオフセット29で配列に追加されます。


最後から2番目の\TT{pop}命令は\ac{TOS}の要素を破棄するので、\TT{return}が実行されたときにスタックが空（またはバランス）になります。